\section{Siegel--Walfisz and Bombieri--Vinogradov}
\label{sec:swbv}

\paragraph{Status of this section.}
The full two uniform distribution theorems are \emph{not proved here}.
We give exact arithmetic targets, prove several finite ingredients, and
locate a substantive missing uniformity step. This qualification does not
apply to the separate Dirichlet proof elsewhere in this report.
The distinction is essential: the usual full Siegel--Walfisz theorem has
ineffective constants, and standard proofs of the full
Bombieri--Vinogradov theorem inherit that issue. Akbary and Hambrook
explain this dependence and prove an explicit restricted-modulus
variant~\cite{AkbaryHambrook}. Johnston's 2026 paper still describes all
known proofs of the unrestricted theorem as ineffective; its new
effective conclusions concern sifting applications~\cite{Johnston2026}.
These facts do not prove nonprovability in iPRA. They identify a proof
obligation which neither a classical quotation nor a finite numerical
check discharges.

\subsection{The targets and their quantifiers}

Throughout this section, $q\geq1$ and $a$ is a reduced residue modulo $q$.
Put
\[
 \Lambda(n)=
 \begin{cases}
  \log p,&n=p^j\text{ for a prime }p\text{ and }j\geq1,\\
  0,&\text{otherwise},
 \end{cases}
 \qquad
 \psi(x;q,a)=\sum_{\substack{1\leq n\leq x\\n\equiv a\pmod q}}\Lambda(n).
\]
For an integer $X\geq1$, define the finite maximum
\[
 E(X;q)=
 \max_{\substack{0\leq y\leq X\\y\in\mathbb N}}
 \max_{\substack{0\leq a<q\\(a,q)=1}}
 \left|\psi(y;q,a)-\frac{y}{\varphi(q)}\right|.
\]
The convention $a=0$ works when $q=1$. Replacing integer $y$ by real
$y$ changes this maximum by at most $1$: on $[n,n+1)$ the sum is constant,
whereas $y/\varphi(q)$ changes by at most $1$.

An exponential Siegel--Walfisz target is
\begin{equation}
 \begin{gathered}
 \forall D\geq1\ \exists c_D,C_D,X_D>0\ \forall x\geq X_D:\\
 \left|\psi(x;q,a)-\frac{x}{\varphi(q)}\right|
 \leq C_Dx\exp(-c_D\sqrt{\log x})
 \quad(q\leq(\log x)^D).
 \end{gathered}
 \label{eq:sw-analytic-target}
\end{equation}
The quantification over $q,a$ in this and similar displays is universal.
An equivalent customary choice of parameters may put ineffectivity in
the multiplicative constant or threshold instead of the exponential
constant. Moving it between parameters does not supply effective
witnesses.

A customary Bombieri--Vinogradov target is
\begin{equation}
 \forall A\geq1\ \exists B,C,X_0\ \forall X\geq X_0:
 \qquad
 \sum_{q\leq \sqrt X/(\log X)^B}E(X;q)
 \leq \frac{CX}{(\log X)^A}.
 \label{eq:bv-analytic-target}
\end{equation}
Even when a formula for $B$ is known, constructive content is still
required for $C$ and $X_0$. A statement with the existential constants
placed after $X$, or a computation for each prescribed finite interval,
does not establish either target.

\subsection{Rational names with a proved error bound}

Here is one deliberately elementary implementation of the logarithms
appearing in these targets. For $n\geq1$ and precision $s\geq0$, set
\[
 z_n=\frac{n-1}{n+1},\qquad M(n,s)=n(n+s+3),\qquad
 L_s(n)=2\sum_{j=0}^{M(n,s)-1}\frac{z_n^{2j+1}}{2j+1}.
\]
These are primitive recursive rational expressions. We use signed
integer numerators and positive integer denominators; every comparison
is implemented by cross multiplication.

\begin{lemma}[Finite logarithm enclosure]
For $r\geq s$,
\[
 0\leq L_r(n)-L_s(n)\leq2^{-s}.
\]
In the ordinary real interpretation,
$0\leq\log n-L_s(n)\leq2^{-s}$.
\end{lemma}
\begin{proof}
If $n=1$, all terms vanish. For $n\geq2$, extending a finite positive
tail to its geometric majorant gives
\[
 L_r(n)-L_s(n)
 \leq \frac{2z_n^{2M(n,s)+1}}{1-z_n^2}.
\]
This inequality follows directly by multiplying the finite geometric
sum by $1-z_n^2$; no infinite sum is needed for the arithmetic assertion.
Now $z_n\leq1-1/n$ and
\[
 \left(1-\frac1n\right)^n\leq\frac12,
\]
since the binomial theorem gives
$(1+1/(n-1))^n\geq1+n/(n-1)>2$.
Also
\[
 \frac{2}{1-z_n^2}=\frac{(n+1)^2}{2n}
 \leq2n\leq2^{n+1}.
\]
Consequently the geometric majorant is at most
\[
 2n\,2^{-2(n+s+3)}
 \leq2^{-n-2s-5}\leq2^{-s}.
\]
The last estimate is a finite integer inequality. The external
identification with $\log n$ follows by integrating the geometric
identity for $2/(1-z^2)$, or by the usual power series for
$\log((1+z)/(1-z))$. Internally the sequence $L_s(n)$ and the proved
finite difference bound are the logarithm name; its external
interpretation is not an additional arithmetic axiom.
\end{proof}

Set $\Lambda_s(n)=L_s(p)$ if $n=p^j$ is a prime power, and zero
otherwise. Bounded trial division decides that case. Let
$\Psi_s(y;q,a)$ be the corresponding finite progression sum, and let
$E_s(X;q)$ be the finite maximum formed with $\Psi_s$ instead of $\psi$.
For $r\geq s$ and $y\leq X$, the preceding lemma gives
\begin{align}
 0&\leq\Psi_r(y;q,a)-\Psi_s(y;q,a)\leq X2^{-s},
 \label{eq:psi-finite-error}\\
 |E_r(X;q)-E_s(X;q)|&\leq X2^{-s}.
 \label{eq:max-finite-error}
\end{align}
For the second inequality, use
$\bigl||u|-|v|\bigr|\leq|u-v|$ and the analogous elementary inequality
for finite maxima. Thus the error for a sum of $Q$ such maxima is
\emph{$QX2^{-s}$}, not merely $X2^{-s}$.

Define the integer functions
\[
 t(X)=\lfloor\log_2X\rfloor,\qquad
 h(t,K)=\lfloor\sqrt t\rfloor\mathbin{\mathrm{div}}K,
 \qquad
 Q_B(t)=\left\lfloor\frac{2^{\lfloor t/2\rfloor}}{t^B}\right\rfloor.
\]
The logarithm notation in $t(X)$ is only descriptive: its implementation
is bounded search for the largest $j\leq X$ with $2^j\leq X$.
Integer square root is likewise bounded search.

One positive arithmetic form of the exponential SW target is
\begin{equation}
\begin{split}
 &\forall D\geq1\ \exists K,C,T\geq1\ \forall X\geq2^T
       \ \forall s,q,a:\\[-2pt]
 &\quad
 1\leq q\leq t(X)^D,\quad a<q,\quad(a,q)=1
 \ \Longrightarrow\quad
 \left|\Psi_s(X;q,a)-\frac{X}{\varphi(q)}\right|
 \leq \frac{CX}{2^{h(t(X),K)}}+\frac{X}{2^s}.
\end{split}
\label{eq:sw-arithmetic-target}
\end{equation}
All variables here are natural numbers, and all displayed fractions are
rational codes. Clearing denominators makes the matrix decidable.
The scale $t(X)$ differs from $\log X$ by a constant factor; replacing
$D$ by $D+1$ and increasing an explicit threshold interconverts the
modulus ranges. The integer $h(t,K)$ differs from $\sqrt t/K$ by less
than $2$. These observations explain the correspondence with
\eqref{eq:sw-analytic-target}, including its exponential error.

A dyadic positive arithmetic BV target is
\begin{equation}
\begin{split}
 &\forall A\geq1\ \exists B\geq0,\ C,T\geq1
       \ \forall t\geq T\ \forall s\geq0:\\
 &\qquad
 \sum_{1\leq q\leq Q_B(t)}E_s(2^t;q)
 \leq \frac{C2^t}{t^A}+Q_B(t)2^{t-s}.
\end{split}
\label{eq:bv-arithmetic-target}
\end{equation}
Negative exponents in $2^{t-s}$ denote rational reciprocals, not integer
subtraction. One may instead multiply through by $2^s t^A$.
The universal precision quantifier encodes a closed inequality between
Cauchy names. It does not request an unbounded search for a convenient
precision. The dyadic presentation recovers arbitrary integer endpoints
by taking the next power of two, using the maximum over $y$, and changing
the constants and logarithmic exponents if necessary.

\subsection{A complete implication from effective BV to a uniform SW input}

The following reduction shows why one cannot remove the constructive
issue merely by appealing to averaging. Its conclusion is the
\emph{arbitrary logarithmic-power saving} form of SW, rather than the
stronger single exponential estimate.

\begin{proposition}[Witness-preserving BV reduction]
Suppose witnesses $B(A),C(A),T(A)$ for
\eqref{eq:bv-arithmetic-target} have been supplied. For any $D\geq1$ put
\[
 T'(A,D)=\max\{T(A),\,16(D+B(A)+1)^2,\,4\}.
\]
Then for $X\geq2^{T'(A,D)}$, $q\leq t(X)^D$, and $(a,q)=1$,
\[
 \left|\Psi_r(X;q,a)-\frac{X}{\varphi(q)}\right|
 \leq \frac{2C(A)X}{t(X)^A}+X2^{-r}
 \qquad\text{for every }r\geq0.
\]
All transformations of the supplied witnesses are primitive recursive.
\end{proposition}
\begin{proof}
First, for integers $k\geq0$ and $t\geq16(k+1)^2$, we prove
\begin{equation}
 t^k\leq2^{\lfloor t/2\rfloor}.
 \label{eq:integer-power-inclusion}
\end{equation}
Put $m=\lfloor\sqrt t\rfloor$. Then $m\geq4(k+1)$ and
$t<(m+1)^2\leq2^{m+1}$. To check the last inequality, use $4^2=2^4$
as the base case and $(u+1)^2\leq2u^2$ for $u\geq3$ as the induction
step. Moreover $2k(m+1)\leq m^2\leq t$. Therefore
$t^k\leq2^{k(m+1)}\leq2^{\lfloor t/2\rfloor}$, as required.

Set $k=D+B(A)$. If $q\leq t^D$, then
$qt^{B(A)}\leq t^k\leq2^{\lfloor t/2\rfloor}$, so
$q\leq Q_{B(A)}(t)$. Since each summand in the BV bound is nonnegative,
every individual $E_s(2^t;q)$ is bounded by that sum.

For the stated $X$, let $t=\lceil\log_2X\rceil$ and
$\ell=t(X)$. Thus $X\leq2^t<2X$, $\ell\leq t$, and
$q\leq\ell^D\leq t^D$. For any $s\geq r$,
\eqref{eq:psi-finite-error} and the BV inequality give
\[
 \left|\Psi_r(X;q,a)-\frac{X}{\varphi(q)}\right|
 \leq X2^{-r}+\frac{C(A)2^t}{t^A}
                 +Q_{B(A)}(t)2^{t-s}.
\]
We may remove the last term because this rational inequality holds for
every $s\geq r$. Explicitly, if rational $u,v,c\geq0$ satisfy
$u\leq v+c2^{-s}$ for every sufficiently large $s$, then $u\leq v$.
If $u>v$, write $u-v=m/n>0$ and $c=a/b$ with positive denominators.
Choose $s\geq r$ with $2^s>an$; then
$c2^{-s}<1/n\leq u-v$, a contradiction. Such an $s$ is given by a
primitive recursive expression, for example $\max(r,an+1)$.
The case split is on decidable rational order. Finally,
$C(A)2^t/t^A\leq2C(A)X/\ell^A$.
\end{proof}

This argument supplies no witnesses for BV. It proves that supplying
them would also supply the uniform small-modulus power-saving estimates
needed in standard analytic treatments. In particular, the aggregate
precision allowance cannot be dropped when a single summand is taken.

\subsection{The precise exceptional-zero obstruction in the usual route}

Suppose a separately established explicit formula has given
\begin{equation}
 \left|\psi(x;q,a)-\frac{x}{\varphi(q)}
        +\frac{\chi_*(a)x^\beta}{\beta\varphi(q)}\right|
 \leq Kx\exp(-\kappa\sqrt{\log x}),
 \label{eq:conditional-exceptional-formula}
\end{equation}
with the correction absent unless an exceptional primitive real
character induces a character modulo $q$. This estimate, with its
quantitative hypotheses, is an \emph{unproved analytic input in this
section}. Assume $\beta\geq1/2$ and $1-\beta\geq\delta>0$.
The fully elementary consequence is
\[
 \left|\psi(x;q,a)-\frac{x}{\varphi(q)}\right|
 \leq Kx e^{-\kappa\sqrt{\log x}}
             +\frac{2x e^{-\delta\log x}}{\varphi(q)}.
\]
If $q\leq(\log x)^D$ and an additional estimate supplied an explicit
$c_D>0$ such that $1-\beta\geq c_Dq^{-1/(2D)}$, then
\[
 (1-\beta)\log x\geq c_D\sqrt{\log x}.
\]
The desired SW conclusion would follow with multiplier $K+2$ and
exponential constant $\min(\kappa,c_D)$. The calculation locates the
missing effective Siegel-type input; it does not establish that input.

By contrast, a lower bound proportional to $q^{-1/2}$ produces only
$(\log x)^{1-D/2}$ in the exponent at the endpoint
$q=(\log x)^D$. This does not even grow when $D\geq2$.
The still weaker bound with an extra logarithmic denominator cannot fix
that defect. An exceptional primitive character also induces characters
at multiples of its conductor. Removing one modulus therefore does not
remove all exceptional contributions. Effective variants excluding
\emph{every multiple} of an exceptional conductor are discussed by
Tao~\cite{TaoBV2015,TaoANT2014}; they are different statements from the
unrestricted targets above.

\subsection{A complete rational-weight Vaughan identity}

We now prove a useful ingredient which has no analytic uniformity gap.
For a numerical bound $N$, supply a finite list of rational weights
$w(p)$ for primes $p\leq N$, encoded by one integer. All assertions
below concern $n\leq N$; no function quantifier is required. Equivalently,
one may fix any primitive recursive weight function. Define
\[
 L_w(n)=\sum_{p\mid n}v_p(n)w(p),\qquad
 \Lambda_w(n)=
 \begin{cases}
 w(p),&n=p^j,\ j\geq1,\\
 0,&\text{otherwise}.
 \end{cases}
\]
Let $\mu$ be the M\"obius function, $\mathbf1(n)=1$ for $n\geq1$, and
$\varepsilon(n)=1$ for $n=1$ and zero otherwise. For functions on the
positive integers, write
\[
 (f*g)(n)=\sum_{d\mid n}f(d)g(n/d),\qquad
 f_{\leq U}(n)=f(n)\mathbf1_{n\leq U},\quad f_{>U}=f-f_{\leq U}.
\]
Every operation at argument $n$ is a bounded finite computation.

\begin{proposition}[Weighted Vaughan identity]
For positive integer cutoffs $U,V$,
\begin{align*}
 \Lambda_w={}&(\Lambda_w)_{\leq V}
        +\mu_{\leq U}*L_w
        -\mu_{\leq U}*\mathbf1*(\Lambda_w)_{\leq V}\\
       &+\mu_{>U}*\mathbf1*(\Lambda_w)_{>V}.
\end{align*}
This identity holds over rational arithmetic for arbitrary $w$.
\end{proposition}
\begin{proof}
Two finite divisor identities suffice:
\[
 \mu*\mathbf1=\varepsilon,\qquad
 \mathbf1*\Lambda_w=L_w.
\]
For the first, the squarefree divisors of $n$ correspond to subsets of
its distinct prime factors, and their M\"obius values sum to
$(1-1)^r$. This is zero unless $n=1$. For the second, the prime powers
$p,p^2,\ldots,p^{v_p(n)}$ contribute exactly $v_p(n)w(p)$.
Associativity and commutativity of convolution follow by reindexing
finite factorizations $abc=n$. Hence $\mu*L_w=\Lambda_w$, and
\begin{align*}
 \Lambda_w
 &=\mu_{\leq U}*L_w+\mu_{>U}*\mathbf1*\Lambda_w\\
 &=\mu_{\leq U}*L_w
   +\mu_{>U}*\mathbf1*(\Lambda_w)_{\leq V}
   +\mu_{>U}*\mathbf1*(\Lambda_w)_{>V}\\
 &=\mu_{\leq U}*L_w
   +(\varepsilon-\mu_{\leq U}*\mathbf1)*(\Lambda_w)_{\leq V}
   +\mu_{>U}*\mathbf1*(\Lambda_w)_{>V}.
\end{align*}
Expanding the middle term proves the formula.
\end{proof}

The usual logarithmic identity is the specialization $w(p)=\log p$;
the finite decomposition appears, for example, in Tao's Lemma
18~\cite{TaoBV2015}. For arithmetic implementation take $w(p)=L_s(p)$.
The identity then remains \emph{exact}. One must not replace $L_w(n)$
by the independently computed $L_s(n)$ and assume exact logarithmic
additivity. With the definition above,
\[
 |\log n-L_w(n)|
 \leq \left(\sum_{p\mid n}v_p(n)\right)2^{-s}\leq n2^{-s},
\]
and the approximation allowance is explicit.

\subsection{A finite Gram bound and a large-sieve consequence}

The following version is formulated entirely with Gaussian rationals.
It is a direct certificate theorem suitable for arithmetic coding.

\begin{proposition}[Rational Gram certificate]
Let $A$ be an $R\times N$ matrix over $\mathbb Q(i)$ and let
$K=AA^*$. Suppose symmetric nonnegative rational numbers $k_{rs}$
and a rational $L\geq0$ satisfy
\[
 |K_{rs}|^2\leq k_{rs}^2,\qquad \sum_s k_{rs}\leq L.
\]
Then, for every $a\in\mathbb Q(i)^N$,
\[
 \sum_r\left|\sum_n A_{rn}a_n\right|^2
       \leq L\sum_n|a_n|^2.
\]
\end{proposition}
\begin{proof}
For $k>0$ and $|z|^2\leq k^2$, expanding the nonnegative square
$|ku-zv|^2$ proves
\[
 2\operatorname{Re}(\overline u zv)
       \leq k(|u|^2+|v|^2).
\]
Indeed multiply the desired inequality by $k$, and bound
$|z|^2|v|^2$ by $k^2|v|^2$. For $k=0$, the hypothesis forces $z=0$,
so the conclusion still holds. Applying this inequality to each pair
of Gram entries and using symmetry gives
\[
 \|A^*b\|^2
   =\sum_{r,s}\overline b_rK_{rs}b_s
   \leq\sum_r|b_r|^2\sum_s k_{rs}
   \leq L\|b\|^2.
\]
Set $b=Aa$. Finite Cauchy--Schwarz now yields
\[
 \|b\|^4=|\langle a,A^*b\rangle|^2
       \leq\|a\|^2\|A^*b\|^2
       \leq L\|a\|^2\|b\|^2.
\]
If $\|b\|^2=0$, there is nothing to prove; otherwise divide by that
positive rational. The needed Cauchy--Schwarz inequality has the exact
sum-of-squares proof
\[
 \left(\sum_i|u_i|^2\right)\left(\sum_i|v_i|^2\right)
  -\left|\sum_i\overline u_i v_i\right|^2
   =\sum_{i<j}|u_iv_j-u_jv_i|^2\geq0.
\]
Thus the argument uses only finite rational identities, inequalities,
and a decidable zero test.
\end{proof}

For clarity we also record the associated elementary Fourier estimate,
while distinguishing its analytic notation from the rational theorem.
Let $\alpha_1,\ldots,\alpha_R$ be rational points of
$\mathbb R/\mathbb Z$ separated by a positive rational $\delta$, let
$a_n$ be Gaussian-rational coefficients, and let $e(t)=\exp(2\pi it)$.
If $R\geq2$, then
\begin{equation}
 \sum_{r=1}^R\left|
  \sum_{n=M+1}^{M+N}a_ne(n\alpha_r)\right|^2
 \leq\left(N+\delta^{-1}H_{R-1}\right)\sum_n|a_n|^2,
 \qquad H_j=\sum_{v=1}^j\frac1v.
 \label{eq:large-sieve-harmonic}
\end{equation}
For $R=1$, ordinary Cauchy--Schwarz gives the bound with coefficient $N$.

Here is the complete elementary human argument for
\eqref{eq:large-sieve-harmonic}. The diagonal Gram entries are $N$.
For distinct frequencies the finite geometric identity gives
\[
 |K_{rs}|
 \leq\frac{2}{|1-e(\alpha_r-\alpha_s)|}
 =\frac1{|\sin\pi(\alpha_r-\alpha_s)|}
 \leq\frac1{2\|\alpha_r-\alpha_s\|}.
\]
The last step uses $\sin(\pi u)\geq2u$ for $0\leq u\leq1/2$,
the chord inequality from concavity of sine on $[0,\pi/2]$.
Relative to a fixed frequency, divide the other frequencies into two
semicircles. The $j$th closest point on either side has distance at least
$j\delta$. If those sides contain $m$ and $n$ points, their row
contribution is bounded by
\[
 \frac{H_m+H_n}{2\delta}\leq\frac{H_{R-1}}\delta.
\]
The Gram proof applies and gives the stated result. A literal iPRA
translation additionally needs the elementary coded sine inequality
and certified approximations to the matrix entries; importing a theorem
about classical real sine into Lean would not constitute that
translation.

For distinct reduced fractions $a/q$ with $q\leq Q$, circular
separation is at least $Q^{-2}$ and there are at most $Q^2$ points.
Thus the weaker, logarithmically inflated large-sieve bound
\[
 \sum_{q\leq Q}\ \sum_{\substack{0\leq a<q\\(a,q)=1}}
   \left|\sum_n a_ne(an/q)\right|^2
 \leq \left(N+Q^2H_{Q^2}\right)\sum_n|a_n|^2
\]
follows. This establishes a useful mean-square ingredient. It does not
establish the character-conductor bookkeeping, Type I and Type II
bounds, or the uniform small-conductor estimates required by full BV.

\subsection{Published effective substitutes and formalization boundary}

A safe explicit alternative is Akbary--Hambrook's Theorem 1.3.
For $x\geq4$ and $1\leq Q_1\leq Q\leq\sqrt x$, it bounds the
average discrepancy from $\psi(y)/\varphi(q)$, over moduli
$q\leq Q$ whose \emph{least prime factor exceeds $Q_1$}, by
\[
 c_1\left(
 \frac{4x}{Q_1}+4\sqrt x\,Q+18x^{2/3}\sqrt Q
       +5x^{5/6}\log\frac{eQ}{Q_1}
 \right)(\log x)^{9/2},
\]
where the paper specifies $c_1$ explicitly~\cite{AkbaryHambrook}.
We cite this as a further formalization target, not as an iPRA theorem
proved in the present package. The least-prime-factor restriction must
not be weakened to merely $q>Q_1$. Similar caution applies to
exceptional-conductor formulations: a constructive implementation
needs usable exceptional data or zero-free certificates, rather than a
nonconstructive choice of whether a boundary zero exists.

The Lean work for the proved finite results can be organized as follows.
First implement bounded divisor enumeration, prime-power recognition,
M\"obius values, valuations, rational codes, Gaussian-rational matrices,
finite sums, and maxima. Prove the two divisor identities and reindexing
lemmas; the weighted Vaughan formula then follows by the displayed
algebra. Implement $L_s$, prove its finite geometric bound, and lift it
to $\Psi_s$ and $E_s$. Encode
\eqref{eq:sw-arithmetic-target} and
\eqref{eq:bv-arithmetic-target} with their quantifier order unchanged.
The integer lemma \eqref{eq:integer-power-inclusion}, the rational-gap
argument, and the BV reduction are then finite arithmetic proofs.
Finally implement the Gram certificate, using the displayed
sum-of-squares identity in place of any infinite-dimensional
functional analysis.

For the original SW and BV conclusions, the outstanding work is
mathematical: quantitative character theory, explicit formulas,
zero-free estimates with verified moduli, complete mean-value bounds,
and effective witnesses for the required uniform Siegel-type input.
Their absence is recorded as an absence of proof, not disguised as
ordinary Lean placeholders or inferred to be a theorem of
nonprovability.
